% !PRINTED 463
\weberpage{463}
%<<CONT
of degree $17$, which does not exist, and consequently $K'$ is
identical with $aK$.

\begin{center}
\gesperrt{There is therefore, apart from a constant factor, only
one invariant of degree $21$.}
\end{center}

But now the system of the 21 principal axes is also an invariant
of degree $21$, and from this it is to be concluded:

\begin{center}
\gesperrt{The invariant $K$ splits into 21 linear factors which,
set equal to zero, represent the principal axes of the group.}
\end{center}

We can then form an invariant of degree $42$,
namely
\weq{2}{\Delta^7 - kC^3 = Q,}
in which $k$ is an arbitrary constant, and this constant can
be determined so that the curve $Q$ passes through an arbitrarily given
point on $f$, and it must then also pass through all the points
connected with this one. Thus we can determine $k$ in $Q$ so
that the curve $Q$ cuts out of the curve $f$ an arbitrarily
given system of connected points.

If now $\Phi$ is an arbitrary invariant not divisible by $f$,
which passes $h_1, h_2, h_3$ times through the poles $P_7, P_3, \overline{P}_2$,
and moreover through arbitrary systems $S_1, S_2, \ldots$
of connected points on $f$, which may also partly coincide,
then we first form, according to (2), the forms
$Q_1, Q_2, \ldots$, which vanish in the systems $S_1, S_2, \ldots$,
and then the form
\[
\Psi = \Phi - a\Delta^{h_1} C^{h_2} K^{h_3} Q_1 Q_2 \ldots
\]

The curve $\Psi$ passes, for every value of the constant $a$, through
all the points of intersection of $\Phi$ with $f$, and if we therefore
determine $a$ so that $\Psi$ passes through some point of $f$
different from these, then $\Psi$ is divisible by $f$, hence
\[
\Phi = a\Delta^{h_1} C^{h_2} K^{h_3} Q_1 Q_2 \cdots + f\Phi_1.
\]
In this, $\Phi_1$ is now again an invariant, but of lower
degree than $\Phi$, and by complete induction the
theorem is hereby proved:

\begin{center}
\gesperrt{Every invariant of the group can be represented as an
integral rational function of the four fundamental invariants
$f, \Delta, C, K$.}
\end{center}

If one determines the constant $k$ in (2) so that the curve $Q$ passes
through one of the poles $\overline{P}_2$, then it can pass through no point not
%CONT>>
